% GENERATED FILE. Edit canonical lessons, then run npm run generate:books.
% canonical-source-hash: fnv1a64:c34d0ad6f052127d
% canonical-lessons: KA-C487-candamaruta, KA-C487-suntaragali, KA-C487-abbarisu, KA-C487-bhukusita, KA-C487-sunami

\chapter{Storms and Disasters}
\label{ch:ka-storms-and-disasters}
\hlchaptermodality{487}

\begin{chapteropening}
\textbf{By the end of this chapter:} \emph{I can say a cyclone, a whirlwind, to rage, a landslide and a tsunami.}

Complete the last lesson of chapter 487: I can say a cyclone, a whirlwind, to rage, a landslide and a tsunami.
\end{chapteropening}

\section[\texorpdfstring{\kn{ಚಂಡಮಾರುತ} (caṇḍamāruta)}{ಚಂಡಮಾರುತ (caṇḍamāruta)}]{\kn{ಚಂಡಮಾರುತ} (caṇḍamāruta) — a cyclone}

\label{lesson:KA-C487-candamaruta}



\begin{quote}
Before the new one: say the Kannada for to shroud, then the Kannada for snowfall.
\end{quote}

\subsection*{You'll want to know: \kn{ಚಂಡಮಾರುತ}}
\textbf{\kn{ಚಂಡಮಾರುತ}} — \emph{caṇḍamāruta} — \textquotedblleft{}a cyclone\textquotedblright{}.

\begin{quote}
\kn{ಚಂಡಮಾರುತದಿಂದ} \kn{ಮೀನುಗಾರರು} \kn{ಸಮುದ್ರಕ್ಕೆ} \kn{ಹೋಗಲಿಲ್ಲ}.
\end{quote}

\begin{quote}
\emph{caṇḍamārutadinda mīnugāraru samudrakke hōgalilla.}
\end{quote}

\begin{quote}
Because of the cyclone the fishermen did not go to sea.
\end{quote}

\kn{ಚಂಡ}, fierce, plus \kn{ಮಾರುತ}, wind.

\begin{cousinweb}
\textbf{\kn{ಚಂಡಮಾರುತ}} (\emph{caṇḍamāruta}) is built from two Sanskrit words: \emph{caṇḍa}, \textquotedblleft{}passionate, fierce\textquotedblright{} (CDIAL 4584), and \emph{māruta}, \textquotedblleft{}windy; the wind\textquotedblright{} (CDIAL 10068).
\end{cousinweb}

\begin{grammarlens}[title={What you've built}]
One more word for getting around.
\end{grammarlens}

\subsection*{Your turn}
\begin{itemize}
\raggedright
  \item \emph{Say it:} \emph{caṇḍamāruta}
  \item \emph{Say it:} \emph{caṇḍamāruta}, once more
  \item \emph{Say it:} \emph{caṇḍamārutadinda mīnugāraru samudrakke hōgalilla.} — Because of the cyclone the fishermen did not go to sea.
  \item \emph{Recall:} say the Kannada for to shroud, then the Kannada for snowfall, then say \emph{caṇḍamāruta} again
\end{itemize}

\begin{tcolorbox}[breakable,colback=teal!4,colframe=teal!35!black,title={Before you move on}]
What does \kn{ಚಂಡಮಾರುತ} mean? (\textquotedblleft{}A cyclone\textquotedblright{}.) Say it once more.
\end{tcolorbox}
\section[\texorpdfstring{\kn{ಸುಂಟರಗಾಳಿ} (suṇṭaragāḷi)}{ಸುಂಟರಗಾಳಿ (suṇṭaragāḷi)}]{\kn{ಸುಂಟರಗಾಳಿ} (suṇṭaragāḷi) — a whirlwind}

\label{lesson:KA-C487-suntaragali}



\begin{quote}
Before the new one: say the Kannada for snowfall, then the Kannada for a cyclone.
\end{quote}

\subsection*{You'll want to know: \kn{ಸುಂಟರಗಾಳಿ}}
\textbf{\kn{ಸುಂಟರಗಾಳಿ}} — \emph{suṇṭaragāḷi} — \textquotedblleft{}a whirlwind\textquotedblright{}.

\begin{quote}
\kn{ಸುಂಟರಗಾಳಿ} \kn{ಹೊಲದ} \kn{ಧೂಳನ್ನು} \kn{ಮೇಲಕ್ಕೆ} \kn{ಎತ್ತಿತು}.
\end{quote}

\begin{quote}
\emph{suṇṭaragāḷi holada dhūḷannu mēlakke ettitu.}
\end{quote}

\begin{quote}
A whirlwind lifted the dust of the field up into the air.
\end{quote}

A spinning wind; \kn{ಬಿರುಗಾಳಿ} is a storm.

\begin{cousinweb}
\textbf{\kn{ಸುಂಟರಗಾಳಿ}} (\emph{suṇṭaragāḷi}) ends in native \kn{ಗಾಳಿ}, \textquotedblleft{}wind, air\textquotedblright{} (DEDR 1499); the source of its first part, \emph{suṇṭara}, is not settled.
\end{cousinweb}

\begin{grammarlens}[title={What you've built}]
One more word for getting around.
\end{grammarlens}

\subsection*{Your turn}
\begin{itemize}
\raggedright
  \item \emph{Say it:} \emph{suṇṭaragāḷi}
  \item \emph{Say it:} \emph{suṇṭaragāḷi}, once more
  \item \emph{Say it:} \emph{suṇṭaragāḷi holada dhūḷannu mēlakke ettitu.} — A whirlwind lifted the dust of the field up into the air.
  \item \emph{Recall:} say the Kannada for snowfall, then the Kannada for a cyclone, then say \emph{suṇṭaragāḷi} again
\end{itemize}

\begin{tcolorbox}[breakable,colback=teal!4,colframe=teal!35!black,title={Before you move on}]
What does \kn{ಸುಂಟರಗಾಳಿ} mean? (\textquotedblleft{}A whirlwind\textquotedblright{}.) Say it once more.
\end{tcolorbox}
\section[\texorpdfstring{\kn{ಅಬ್ಬರಿಸು} (abbarisu)}{ಅಬ್ಬರಿಸು (abbarisu)}]{\kn{ಅಬ್ಬರಿಸು} (abbarisu) — to rage; to roar}

\label{lesson:KA-C487-abbarisu}



\begin{quote}
Before the new one: say the Kannada for a cyclone, then the Kannada for a whirlwind.
\end{quote}

\subsection*{You'll want to know: \kn{ಅಬ್ಬರಿಸು}}
\textbf{\kn{ಅಬ್ಬರಿಸು}} — \emph{abbarisu} — \textquotedblleft{}to rage; to roar\textquotedblright{}.

\begin{quote}
\kn{ಬಿರುಗಾಳಿಯಲ್ಲಿ} \kn{ಸಮುದ್ರ} \kn{ಅಬ್ಬರಿಸಿತು}.
\end{quote}

\begin{quote}
\emph{birugāḷiyalli samudra abbarisitu.}
\end{quote}

\begin{quote}
In the storm the sea raged.
\end{quote}

Past (it …): \kn{ಅಬ್ಬರಿಸಿತು}, \emph{abbarisitu}.

\kn{ಅಬ್ಬರ} is an uproar: of the sea, of storms and of loud speakers.

\begin{cousinweb}
\textbf{\kn{ಅಬ್ಬರಿಸು}} (\emph{abbarisu}) is native Dravidian: DEDR 367 lists Kannada \emph{abbara}, \textquotedblleft{}a loud cry, a noise\textquotedblright{}, and \emph{abbarisu}, \textquotedblleft{}to cry aloud, to howl\textquotedblright{}.
\end{cousinweb}

\begin{grammarlens}[title={What you've built}]
One more word for getting around.
\end{grammarlens}

\subsection*{Your turn}
\begin{itemize}
\raggedright
  \item \emph{Say it:} \emph{abbarisu}
  \item \emph{Say it:} \emph{abbarisu}, once more
  \item \emph{Say it:} \emph{birugāḷiyalli samudra abbarisitu.} — In the storm the sea raged.
  \item \emph{Recall:} say the Kannada for a cyclone, then the Kannada for a whirlwind, then say \emph{abbarisu} again
\end{itemize}

\begin{tcolorbox}[breakable,colback=teal!4,colframe=teal!35!black,title={Before you move on}]
What does \kn{ಅಬ್ಬರಿಸು} mean? (\textquotedblleft{}To rage; to roar\textquotedblright{}.) Say it once more.
\end{tcolorbox}
\section[\texorpdfstring{\kn{ಭೂಕುಸಿತ} (bhūkusita)}{ಭೂಕುಸಿತ (bhūkusita)}]{\kn{ಭೂಕುಸಿತ} (bhūkusita) — a landslide}

\label{lesson:KA-C487-bhukusita}



\begin{quote}
Before the new one: say the Kannada for a whirlwind, then the Kannada for to rage.
\end{quote}

\subsection*{You'll want to know: \kn{ಭೂಕುಸಿತ}}
\textbf{\kn{ಭೂಕುಸಿತ}} — \emph{bhūkusita} — \textquotedblleft{}a landslide\textquotedblright{}.

\begin{quote}
\kn{ಘಾಟಿಯಲ್ಲಿ} \kn{ಭೂಕುಸಿತದಿಂದ} \kn{ರಸ್ತೆ} \kn{ಮುಚ್ಚಿತು}.
\end{quote}

\begin{quote}
\emph{ghāṭiyalli bhūkusitadinda raste muccitu.}
\end{quote}

\begin{quote}
The ghat road was closed by a landslide.
\end{quote}

\kn{ಭೂ}, earth, plus \kn{ಕುಸಿತ}, a sinking or collapse.

\begin{cousinweb}
\textbf{\kn{ಭೂಕುಸಿತ}} (\emph{bhūkusita}) joins Sanskrit \emph{bhū}, \textquotedblleft{}the earth, the ground\textquotedblright{} (CDIAL 9546), and native \emph{kusita}, \textquotedblleft{}a collapse\textquotedblright{}, from \emph{kusi}, \textquotedblleft{}to sink, to give way\textquotedblright{} (DEDR 1636).
\end{cousinweb}

\begin{grammarlens}[title={What you've built}]
One more word for getting around.
\end{grammarlens}

\subsection*{Your turn}
\begin{itemize}
\raggedright
  \item \emph{Say it:} \emph{bhūkusita}
  \item \emph{Say it:} \emph{bhūkusita}, once more
  \item \emph{Say it:} \emph{ghāṭiyalli bhūkusitadinda raste muccitu.} — The ghat road was closed by a landslide.
  \item \emph{Recall:} say the Kannada for a whirlwind, then the Kannada for to rage, then say \emph{bhūkusita} again
\end{itemize}

\begin{tcolorbox}[breakable,colback=teal!4,colframe=teal!35!black,title={Before you move on}]
What does \kn{ಭೂಕುಸಿತ} mean? (\textquotedblleft{}A landslide\textquotedblright{}.) Say it once more.
\end{tcolorbox}
\section[\texorpdfstring{\kn{ಸುನಾಮಿ} (sunāmi)}{ಸುನಾಮಿ (sunāmi)}]{\kn{ಸುನಾಮಿ} (sunāmi) — a tsunami}

\label{lesson:KA-C487-sunami}



\begin{quote}
Before the new one: say the Kannada for to rage, then the Kannada for a landslide.
\end{quote}

\subsection*{You'll want to know: \kn{ಸುನಾಮಿ}}
\textbf{\kn{ಸುನಾಮಿ}} — \emph{sunāmi} — \textquotedblleft{}a tsunami\textquotedblright{}.

\begin{quote}
\kn{ಸುನಾಮಿಯಿಂದ} \kn{ಕರಾವಳಿಯ} \kn{ಹಳ್ಳಿಗಳು} \kn{ನಾಶವಾದವು}.
\end{quote}

\begin{quote}
\emph{sunāmiyinda karāvaḷiya haḷḷigaḷu nāśavādavu.}
\end{quote}

\begin{quote}
Coastal villages were destroyed by the tsunami.
\end{quote}

A great sea wave raised by an earthquake under the sea.

\begin{cousinweb}
\textbf{\kn{ಸುನಾಮಿ}} (\emph{sunāmi}) is a loan from English \emph{tsunami}, which English took from Japanese.
\end{cousinweb}

\begin{grammarlens}[title={What you've built}]
One more word for getting around.
\end{grammarlens}

\subsection*{Your turn}
\begin{itemize}
\raggedright
  \item \emph{Say it:} \emph{sunāmi}
  \item \emph{Say it:} \emph{sunāmi}, once more
  \item \emph{Say it:} \emph{sunāmiyinda karāvaḷiya haḷḷigaḷu nāśavādavu.} — Coastal villages were destroyed by the tsunami.
  \item \emph{Recall:} say the Kannada for to rage, then the Kannada for a landslide, then say \emph{sunāmi} again
\end{itemize}

\begin{tcolorbox}[breakable,colback=teal!4,colframe=teal!35!black,title={Before you move on}]
What does \kn{ಸುನಾಮಿ} mean? (\textquotedblleft{}A tsunami\textquotedblright{}.) Say it once more.
\end{tcolorbox}
